\subsection{Mechanical}

\subsubsection{Original MIRTE Master Platform}

The project is based on the MIRTE Master V2. The platform consists of a mobile base propelled by four individually driven mecanum wheels and a robotic arm, allowing the robot to navigate through an environment and interact with objects.

The original MIRTE Master hardware consists of three main parts:

\begin{enumerate}
\item \textbf{Mobile base}
\item \textbf{Electronics}
\item \textbf{Robotic arm}
\end{enumerate}

The mobile base utilizes four geared DC motors with Hall encoders, connected to 97 mm mecanum wheels. This allows the robot to move in multiple directions without having to turn first, which is useful for navigating in small indoor spaces. The base also contains various sensors, such as rearward facing ultrasonic sensors, an IMU, a 2D LiDAR and an RGBD (depth) camera mounted on the robotic arm.

The electronics of the original platform are based on an Orange Pi 3B and a Raspberry Pi Pico H. The Orange Pi functions as the main computer, while the Pico is used for lower-level control tasks such as controlling the motors and general interfacing. The MIRTE Master also includes a custom-made PCB to connect the motors, sensors, and power system in a structured manner. The main power source for the MIRTE Master is a 5Ah 12V Lion Parkside drill battery. No significant changes were made to the electronics used for this project.

The robotic arm is mounted on top of the mobile base and has four degrees of freedom, which are driven by four serial bus servo motors of varying strengths. A gripper and RGBD camera are mounted at the end of of the arm for object scanning and manipulation.

\subsubsection{Hardware Components}

The most important original hardware components can be seen in \href{https://matt-rbt.github.io/Lab-Cleanup-Robot-using-the-Mirte-Master-Platform/materials/}{Materials}.

These components provide a useful base platform for a mobile manipulation task. However, for this project, the original configuration needed changes. The robot had to be adapted for a laboratory-cleanup task, where it should be able to move around in a laboratory environment, detect waste on the floor, and eventually collect and move objects into bins.

\subsubsection{Design Goals for the Modified Robot}

The goal of our hardware modifications was to make the MIRTE Master more suitable for collecting and transporting small waste objects. This required changes to the chassis and and gripper design.

The main design goals were:

% - Increase the stiffness and reliability of the chassis.

\begin{itemize}
\item Add space for collecting and separating waste.
\item Mount the camera in a useful position for object recognition.
\item Improve the battery mounting system.
\item Improve cable management.
\item Modify the gripper so it can pick up waste objects more reliably.
\end{itemize}

\subsubsection{Chassis Modifications}

The original MIRTE Master chassis was designed as a general-purpose mobile manipulator base. For our use case, the robot also needed to carry waste bins and support the manipulator during picking tasks. Therefore, the chassis was modified to improve strength and usability.

The chassis consists of six side panels and a top and a bottom plate. In the original model, the top and bottom plates were made out of clear PMMA. While PMMA allows the user to see the status of the electronics clearly by looking at the LED's on the inside, it also has a tendency to shatter easily if bent too much. In early testing of the MIRTE Master it has already shown that the PMMA was not a good choice for any type of loading resulting in shattered PMMA. Therefore a change the material of all plates from PMMA to aluminium with a similair stiffness was essential, the difference of this being that aluminium can bend more and will not shatter. To still be able to see the status of the electronics by the LEDs, the six side panel elements were 3D-printed using a slightly translucent material.

Two separate waste bins were added to the robot. These bins allow the robot to collect, store and seperate objects after picking them up.

\subsubsection{Gripper Modifications}

The gripper is an important part of the laboratory-cleanup robot because it directly interacts with the objects that need to be collected. The original MIRTE Master gripper is useful as a general robotic gripper, but the cleanup task requires it to handle different small objects reliably.

The original gripper geometry is suitable for holding small and large objects, however, the gripping surface is not. Electronic trash tends to have small contact points for the gripper so good gripping performance is a must. For improving the grip on objects the contact patches are layered with a layer of soft 83A TPE, coated with an additional layer of Bison Rubber anti-slip coating.

During the initial testing phase of the RGBD camera, it was found that the mounting location for it was too low to the ground for it to be able to confidently display any accurate results regarding the distance towards objects. While this was the intended way for the camera to be implemented, for this project it was a better plan to mount this RGBD camera onto the end of the arm, right above the gripper. This gives it elevation which makes it able to look slightly more downward and therefore it could more accurately display distances. This also makes the need for the secondary RGB camera mounted onto the gripper insignificant, as the RGBD camera can also be used for object detection.

\subsubsection{Summary of Hardware Changes}

\bigskip\noindent
\begin{tabular}{p{\dimexpr 0.333\linewidth-2\tabcolsep}p{\dimexpr 0.333\linewidth-2\tabcolsep}p{\dimexpr 0.333\linewidth-2\tabcolsep}}
\toprule
Subsystem & Original design & Modified version \\
\hline
Chassis & Standard MIRTE Master frame with PMMA & Strengthened top and bottom plates with aluminium \\
Storage & No dedicated cleanup storage & Two separate waste bins added \\
Camera & RGBD camera mounted at the front of the chassis & RGBD camera mounted on the gripper \\
Wiring & Standard cable routing & Slightly improved cable management \\
Gripper & Original MIRTE gripper & Improved gripping surface \\
\bottomrule
\end{tabular}

\bigskip

\subsubsection{Current Design}

The complete model can be explored by using the interactive 3D interface below.